\chapter{ACTH Stimulation Test}

\section*{What Is This Test?}
The ACTH stimulation test, also called the cosyntropin stimulation test, evaluates the ability of the adrenal cortex to produce cortisol after stimulation with synthetic adrenocorticotropic hormone. It is the standard confirmatory test for suspected adrenal insufficiency.

\section*{Alternative Names}
Cosyntropin Test; Short Synacthen Test; Adrenal Stimulation Test; Cortrosyn Stimulation Test.

\section*{Why This Test Is Ordered}
It is used when primary or central adrenal insufficiency is suspected, when morning cortisol is indeterminate, or when recovery of the hypothalamic-pituitary-adrenal axis is being assessed after prolonged glucocorticoid exposure. It does not by itself distinguish primary from secondary adrenal disease; ACTH, electrolytes, renin, aldosterone, and clinical findings provide that context.

\section*{How This Test Is Performed}
Blood is collected for baseline cortisol, and sometimes ACTH. Synthetic cosyntropin is administered according to the laboratory protocol, commonly 250 micrograms intravenously or intramuscularly. Cortisol is measured again at approximately 30 and 60 minutes. Low stimulated cortisol suggests impaired adrenal reserve.

\section*{Preparation, Risks, and Limitations}
Testing is usually performed in the morning. The ordering clinician must review glucocorticoids and other medicines that alter cortisol measurement or adrenal function. Risks are usually limited to venipuncture discomfort and brief medication-related flushing or nausea. Recent exogenous steroid use, critical illness, pregnancy, and assay-specific cutoffs can affect interpretation.

\section*{Understanding Results}
An adequate cortisol rise generally argues against clinically significant adrenal insufficiency, but the exact cutoff is assay-dependent. A poor response supports adrenal insufficiency. A normal response can occur early in recent central adrenal disease, when the adrenal glands have not yet atrophied. Results should be interpreted with the baseline ACTH and clinical picture.

\section*{References}
Endocrine Society clinical practice guidance on diagnosis and treatment of primary adrenal insufficiency.

\clearpage
\section*{Understanding Results and Follow-up}
The laboratory report should identify the specimen, method, units, reference interval or decision limit, and whether the result is positive, negative, abnormal, or inconclusive. Reference intervals vary with age, sex, pregnancy, specimen type, assay, and laboratory; a result outside a range is not automatically a diagnosis, and a result within a range does not always exclude disease. Discuss unexpected results with the ordering clinician, who can decide whether repeat or confirmatory testing, treatment, or referral is appropriate. Do not change medicines or delay urgent care based on an isolated result.


\section*{Regulatory Status and Authorization Date}
Regulatory status applies to the specific assay, instrument, manufacturer, and intended use rather than to the test name alone. No single approval date can be assigned to this test category. For a commercial assay, verify the current product in the relevant regulator's database: the U.S. Food and Drug Administration (FDA) clearance, approval, or authorization; India's Central Drugs Standard Control Organisation (CDSCO) licence or permission; the European Union In Vitro Diagnostic Regulation (EU IVDR) CE certificate issued through the applicable conformity-assessment route; or the United Kingdom Medicines and Healthcare products Regulatory Agency (MHRA) registration and UKCA/CE status. Laboratory-developed tests may not have an individual FDA, CDSCO, EU IVDR, or MHRA approval date.
\clearpage
